% TOP_PROBLEM: {"id": 22, "title": "Continuum Hypothesis", "category_id": 10, "category": {"id": 10, "name": "set_theory", "display_name": "Set Theory", "description": "Foundations of mathematics, infinite sets, and cardinality.", "slug": "set-theory", "order_index": 10, "created_at": "2026-07-31T15:26:25.671Z"}, "status": "open"}
% FIRST_POSTED: null
% ENTRY_KIND: statement_only
% Imported from: batch01.tex; UnsolvedMath ID 22
\documentclass[11pt]{article}
\usepackage[margin=25mm]{geometry}
\usepackage{fontspec}
\setmainfont{Latin Modern Roman}
\usepackage{amsmath,amssymb,mathtools}
\usepackage{unicode-math}
\setmathfont{Latin Modern Math}
\usepackage{xurl,hyperref}
\hypersetup{colorlinks=true,urlcolor=blue}
\setcounter{secnumdepth}{0}
\setlength{\parindent}{0pt}
\setlength{\parskip}{5pt}
\setlength{\emergencystretch}{3em}
\newcommand{\sref}[2]{\hyperlink{source:#1:#2}{[S#2]}}
\newcommand{\cataloguescope}[1]{\paragraph{Recorded target.}#1}
\begin{document}
% UnsolvedMath ID 22; source catalogue code SET-001
\setcounter{section}{21}
\section{Continuum Hypothesis}\label{22}
{\small\sffamily UnsolvedMath ID 22 · Set Theory}
\subsection{Definitions and mathematical statement}

\paragraph{Shared notation.}$\mathbb N=\{1,2,\ldots\}$, and $\mathbb Z,\mathbb Q,\mathbb R,\mathbb C$ denote the integers, rationals, reals and complex numbers. Any other conventions are specified locally.


For sets $X,Y$, write $|X|=|Y|$ when there is a bijection between them, and $|X|<|Y|$ when $X$ injects into $Y$ but there is no such bijection. Let $\aleph_0=|\mathbb N|$, let $\aleph_1$ be the least uncountable cardinal, and put $2^{\aleph_0}=|\mathcal P(\mathbb N)|=|\mathbb R|$, where $\mathcal P(X)$ is the power set of $X$. Work in Zermelo--Fraenkel set theory with the axiom of choice (ZFC).
\paragraph{Statement recorded in the supplied source.}
The continuum hypothesis is the assertion
\[
 2^{\aleph_0}=\aleph_1,
\]
equivalently, there is no set $X$ with $\aleph_0<|X|<|\mathbb R|$. The historical question asks whether this assertion follows from ZFC or whether its negation follows from ZFC. \sref{22}{2}
Gödel's and Cohen's relative-consistency results settle this decision question by independence: if ZFC is consistent, neither the assertion nor its negation is provable in ZFC. The formal proof in Source~S2 verifies both directions; independence does not assert that the displayed sentence is false. \sref{22}{2}
\subsection{Short English statement}
The continuum hypothesis says that every uncountable subset of the real numbers has the same cardinality as the real numbers. Its truth cannot be decided from ZFC alone, assuming ZFC is consistent.
\subsection{Sources}
\begin{description}
\item[\hypertarget{source:22:1}{S1}] ULAM.AI and the UnsolvedMath Contributors, UnsolvedMath: A Curated Collection of Open Mathematics Problems, record 22 (SET-001). CC BY 4.0. \url{https://huggingface.co/datasets/ulamai/UnsolvedMath}.
\item[\hypertarget{source:22:2}{S2}] Jesse Michael Han and Floris van Doorn, A Formal Proof of the Independence of the Continuum Hypothesis, Proceedings of CPP 2020, 14 pages; arXiv:2102.02901 (2021), Section 1 and the two forcing constructions. \url{https://arxiv.org/pdf/2102.02901}.
\end{description}
\label{end:22}
\end{document}
